\section{Calibration and experiments}\label{sec:q:cal}

\subsection{What has to be identified, and from what}\label{sec:q:cal:blocks}

The model has an unusually favourable identification structure, because the materials block is measured in tonnes and material-flow accounting reports tonnes. Most of the parameters that a conventional growth model would have to infer from prices are here pinned by physical quantities, and the parameters that remain genuinely difficult are few and can be named.

\begin{table}[H]
\centering
\caption{Parameter blocks and their leading identifying moments.}\label{tab:q:cal}
\begin{tabular}{L{3.1cm}L{4.2cm}L{6.4cm}}
\toprule
Block & Parameters & Identified principally by \\
\midrule
Preferences &
$\rho,\eta$ &
long-run real interest rate and the consumption--wealth ratio; $\eta$ additionally restricted by convergence, $\ln\left(1+\rho\right)>\left(1-\eta\right)g$ \\
Damages &
$\kappa,\psi,\varphi$ &
an assumed damage function calibrated to a conventional loss at a reference pollution stock; $\varphi$ governs convexity and is varied \\
Production &
$\beta,\mu_F,\varsigma,\delta$ &
capital and material shares of gross output, the capital--output ratio, the depreciation rate; $\varsigma$ from the long-run price elasticity of material demand \\
Trends &
$A_0,g_A,B_0,g_B$ &
the historical growth rate of output and the historical decline of material intensity $\Omega$, which separates $g_A$ from $g_B$ \\
Material floor &
$\bar R$ &
\emph{not identified by history}; see Section~\ref{sec:q:cal:hard} \\
Recycling &
$\bar a,\xi$ (or $\psi$) &
observed secondary-production shares and yields by material, and the observed cost gradient of recovery at high recycling rates \\
Waste handling &
$\mu,c^c,c_T,\chi_T,d^W$ &
mean residence time of materials in use and in stock; unit costs of collection and treatment; the share of untreated waste reaching the environment \\
Extraction &
$c_N,\mu_N,\kappa_N,\chi_N,S_{\mathrm{ref}}$ &
the historical path of real material prices against cumulative extraction; $\mu_N$ is the single most consequential parameter and is discussed below \\
Exploration &
$c_D,\mu_D,\kappa_D,\chi_D,\bar X_{\max}$ &
exploration expenditure against discoveries; $\bar X_{\max}$ from geological estimates of the ultimately recoverable resource \\
Accounting &
$\omega^Y,\omega^N,\omega^D,\omega^W,\omega^C,\phi^I,\Omega^{N,S},\Omega^{D,S}$ &
economy-wide material flow accounts: waste generation by source, material embodied in construction and durables, and mining overburden ratios \\
Initial stocks &
$K_0,S_0,X_0,P_0,M^K_0,\mathcal W_0$ &
capital stock, reserves and cumulative discoveries, the accumulated pollution stock, in-use material stocks, and accumulated waste stocks, all at the base year \\
\bottomrule
\end{tabular}
\end{table}

Three features of this table carry most of the quantitative content.

\smalltitle{$\Omega$ is directly observable, and it separates the two trends.} The material intensity $\Omega=R/\left(\mathcal D+\phi^IG\right)$ is, up to the composition weights, material input per unit of economic activity --- an object that material flow accounts report and that has a long and well-documented history. In the model its trend is not free: on any balanced path it decays at the rate at which effective material productivity outpaces physical throughput, so the joint history of output growth and intensity decline over-identifies $\left(g_A,g_B\right)$ and provides a specification test. This is the single most useful thing the materials accounting buys for the calibration.

\smalltitle{$\mu_N$ decides the floorless benchmark.} By~\eqref{eq:app:wh:bdp-solution}, the stock-effect elasticity of extraction costs is what selects between sustained growth, a sustained consumption level and managed decline when $\bar R=0$. It is identified by how strongly real material prices have responded to cumulative depletion, holding technology fixed --- which is exactly the quantity that the long historical record of falling real commodity prices makes hard to pin down, since the observed path confounds depletion with technical progress in extraction. The baseline identifies it jointly with $c_N(t)$ from the price and quantity history and reports the result as a range rather than a point.

\smalltitle{The accounting weights matter less than they look.} The relative intensities $\omega^j$ affect the allocation only through $\Omega$ and the wedges $\mathfrak z\Omega^j$, and the wedges are proportional to $\zeta$, which vanishes with $\phi^I$. Their main role is the one Section~\ref{sec:sp:setup:aggregatematerial} assigns them: mapping the model's mass flows onto published accounts. A calibration that gets $\phi^I$ and the total right, and the split across $\omega^j$ roughly, will not move the paths much, and this is worth verifying by perturbation rather than asserting.

\subsection{The two parameters that history cannot identify}\label{sec:q:cal:hard}

\smalltitle{The material floor $\bar R$.} No historical observation identifies $\bar R$, because the floor has never bound: the whole sample is drawn from the interior of the technology, where output responds smoothly to material input. It is the parameter that decides which half of the taxonomy the economy lives in, and it cannot be estimated from a path on which it is slack. The honest response is to treat it as the object the exercise is \emph{about}: the model is solved across a range of $\bar R$ and the results reported as a function of it, with the survival condition $\mu\mathcal B>\bar R$ of~\eqref{eq:app:wh:cbgp:survival} giving the coarse threshold and the exact condition $\mathcal M_\infty>\bar R\mathcal T$ giving the fine one. \emph{The quantitative question is not what $\bar R$ is but how large it would have to be to matter}, and that question the model answers.

\smalltitle{The yield ceiling $\bar a$.} Whether the ceiling is hard or soft is likewise not a property of the observed range. Recovery rates are far from one for every material, so the data speak to the \emph{curvature} $\xi$ --- the cost gradient of pushing recovery higher --- and not to the limit it approaches. Since the taxonomy turns on $\bar a<1$ against $\bar a=1$, this is reported as a case distinction rather than estimated, and the results are presented for both. Section~\ref{sec:app:workhorse:cbgp} narrows what has to be assumed: within the soft case, the tail exponent does not matter for survival, so only the ceiling itself is a live assumption.

\subsection{The aggregation cost of one material}\label{sec:q:cal:aggregation}

The model's single material has to stand for a set of flows whose circularity differs by an order of magnitude: construction minerals, which dominate mass and are recovered at low value; metals, which are recovered at high rates and can in principle be recovered at very high rates; and fossil and biomass flows, much of which are dissipated by design. A single $\left(\bar a,\xi\right)$ cannot represent that heterogeneity, and the direction of the resulting bias is not obvious a priori: a mass-weighted aggregate is dominated by the least circular flows, while a value-weighted aggregate is dominated by the most circular.

Two responses are used. First, the baseline is calibrated on a mass basis, since the ledger is a mass balance and the floor is a mass constraint, and the metals-only calibration is reported alongside it as a bound. Second, the aggregation is treated as a source of parameter uncertainty rather than as a point assumption, and the survival margin is reported as a function of $\left(\bar a,\xi\right)$ over the range the disaggregated evidence supports. A genuinely multi-material version is the natural next model and is not attempted here.

\subsection{Experiments}\label{sec:q:cal:experiments}

The experiments follow the structure of the theory rather than a list of policy scenarios.

\smalltitle{(E1) Where is the economy in the taxonomy?} Solve the planner's problem on the baseline calibration and locate the resulting path in the table of Section~\ref{sec:sp:longrun:taxonomy}: does the loop close, does the material era end, and if so when. Report the retained endowment $\mathcal M_\infty$, the turnover time $\mathcal T$ and the survival ratio $\mathcal M_\infty/\left(\bar R\mathcal T\right)$, which is the single number the theory says matters.

\smalltitle{(E2) How much does circularity buy in the cells where it is not survival?} In the collapse and rest-point cells the circular economy is a transition technology, and Section~\ref{sec:app:workhorse:shutdown} says exactly what it does there: it multiplies the throughput budget without changing the slope of the material price path. Measure the resulting welfare gain, the extension of the material era, and the cumulative pollution avoided, as a function of $\bar a$ and $\xi$.

\smalltitle{(E3) The cost of the three market failures.} Solve the market economy at each corner of the policy space $\left(\phi^W,\phi^z,\phi^P,\phi^X\right)$ and along paths between them. Two questions are separable and both are interesting: the conventional one, how much welfare each instrument is worth, and the one specific to this model, whether any combination of missing instruments moves the economy across the survival margin of Section~\ref{sec:mkt:impl:survival}. A policy failure that costs a few percent of consumption is a familiar object; one that converts perpetual growth into a scheduled shutdown is not, and whether the second is possible on any defensible calibration is the paper's sharpest quantitative question.

\smalltitle{(E4) Property rights over the waste stock.} The dial $\phi^W$ is the only one that is not a tax rate, and its effect is qualitatively different: it does not correct a wedge, it creates the market in which the ledger is priced at all. Report the equilibrium gate fee path under $\phi^W=1$, in particular the date at which it changes sign, since that date is the model's prediction of when waste stops being a disposal problem and becomes a resource --- and the property right is what makes the prediction operative.

\smalltitle{(E5) Sensitivity along the two axes that history cannot pin down.} Report every headline result as a surface over $\left(\bar R,\bar a\right)$, since these are the parameters the taxonomy turns on and the parameters the data cannot deliver.
